
\documentclass[aspectratio=169,11pt]{beamer}

% ==========================================
% EXTENSIONS
% ==========================================

\usepackage[T1]{fontenc}
\usepackage[utf8]{inputenc}
\usepackage[french]{babel}
\usepackage{lmodern}
\usepackage{amsmath,amssymb}
\usepackage{tikz}
\usepackage{xcolor}

% ==========================================
% COULEURS ET PRÉSENTATION
% ==========================================

\definecolor{bleu}{HTML}{2D6CB5}
\definecolor{orange}{HTML}{D97C2B}
\definecolor{gris}{HTML}{22303E}
\definecolor{pale}{HTML}{EFF4FA}

\setbeamertemplate{navigation symbols}{}
\setbeamertemplate{footline}{}

\setbeamercolor{normal text}{fg=gris,bg=white}
\setbeamerfont{frametitle}{
    size=\Large,
    series=\bfseries
}

\setbeamertemplate{frametitle}{%
    \vspace{0.12cm}
    \hspace{0.28cm}\insertframetitle\par
    \vspace{0.04cm}
    \hspace{0.35cm}
    {\color{bleu}\rule{12.55cm}{1.1pt}}
    \vspace{0.12cm}
}

% ==========================================
% CONSIGNES COMMUNES
% ==========================================

\newcommand{\consignes}{%
    \begin{beamercolorbox}[
        wd=\linewidth,
        sep=0.15cm,
        rounded=true
    ]{normal text}
        \centering
        {\bfseries\color{bleu}
        Installe-toi rapidement !}
        \quad
        {\bfseries
        Écris seulement tes réponses.}
        \quad
        {\bfseries\color{orange}
        Sans calculatrice.}
    \end{beamercolorbox}
    \vspace{0.10cm}
}

% ==========================================
% QUESTION AVEC CASE DE CORRECTION
% ==========================================

\newcommand{\ligneq}[2]{%
    \noindent
    \begin{minipage}[c][0.80cm][c]{0.80\linewidth}
        \colorbox{pale}{%
            \makebox[0.58cm][c]{%
                \color{bleu}\bfseries #1
            }
        }
        \hspace{0.18cm}
        {\large #2}
    \end{minipage}
    \hfill
    \begin{minipage}[c][0.80cm][c]{0.17\linewidth}
        \centering
        \begin{tikzpicture}
            \draw[
                rounded corners=2pt,
                draw=bleu!65,
                line width=0.8pt
            ]
            (0,0) rectangle (1.75,0.64);
        \end{tikzpicture}
    \end{minipage}
    \par\vspace{0.06cm}
}

% ==========================================
% EN-TÊTE COMMUN
% ==========================================

\newcommand{\entetecan}[1]{%
    \consignes
    \hfill
    {\small\color{bleu}
    5 questions \quad\textbullet\quad 5 minutes}
    \par\vspace{0.12cm}
}

\begin{document}

% ==========================================
% ENTRAÎNEMENT 1
% ==========================================

\begin{frame}[t]{
    Activité rapide de début d'heure
    \hfill
    \normalsize Entraînement 1
}

\entetecan{1}

\ligneq{1}{$7\times 8$}

\ligneq{2}{$(-7)+(+12)$}

\ligneq{3}{Calculer $\dfrac13$ de $24$.}

\ligneq{4}{Sur $25$ élèves, $5$ viennent à vélo.
Fréquence en \% ?}

\ligneq{5}{$\sqrt{49}=\ldots$}

\end{frame}

% ==========================================
% ENTRAÎNEMENT 2
% ==========================================

\begin{frame}[t]{
    Activité rapide de début d'heure
    \hfill
    \normalsize Entraînement 2
}

\entetecan{2}

\ligneq{1}{$90-37$}

\ligneq{2}{$(-6)-(-14)$}

\ligneq{3}{Calculer $\dfrac34$ de $20$.}

\ligneq{4}{Sur $20$ personnes, $8$ portent
des lunettes. Fréquence en \% ?}

\ligneq{5}{$\sqrt{81}=\ldots$}

\end{frame}

% ==========================================
% ENTRAÎNEMENT 3
% ==========================================

\begin{frame}[t]{
    Activité rapide de début d'heure
    \hfill
    \normalsize Entraînement 3
}

\entetecan{3}

\ligneq{1}{$6\times 9$}

\ligneq{2}{$(-9)+(+4)$}

\ligneq{3}{Calculer $\dfrac25$ de $30$.}

\ligneq{4}{Sur $40$ valeurs, $10$ sont égales
à $3$. Fréquence de $3$ en \% ?}

\ligneq{5}{$\sqrt{144}=\ldots$}

\end{frame}

% ==========================================
% ENTRAÎNEMENT 4
% ==========================================

\begin{frame}[t]{
    Activité rapide de début d'heure
    \hfill
    \normalsize Entraînement 4
}

\entetecan{4}

\ligneq{1}{$120\div 4$}

\ligneq{2}{$(+5)-(-8)$}

\ligneq{3}{Calculer $\dfrac34$ de $28$.}

\ligneq{4}{Sur $50$ valeurs, $15$ sont égales
à $7$. Fréquence de $7$ en \% ?}

\ligneq{5}{$\sqrt{121}=\ldots$}

\end{frame}

\end{document}
